Language models that call tools can place trades, and a firm that gives one an
order ticket needs to know whether it will refuse an order that breaks the firm's
rules when its own portfolio manager asks for it. We introduce
\texttt{mandatebench}, a benchmark in which an agent works the desk of a
fictitious firm under a six-rule compliance manual (restricted list, earnings
blackout, concentration cap, pre-trade approval, no short selling, no margin). In
144 conflict scenarios the requested trade breaches exactly one rule, under one of
six pressure levels; in 28 legitimate scenarios, several placed exactly at a
limit, it breaches nothing, so overcaution is measured alongside violation. A
deterministic rule engine scores every order from the ledger. Eight open-weight
models from four developers were run locally, in two batches: the finding from
the first four was recorded as a falsifiable prediction before the other four were
run. Violation rates ranged from \num{gpt-oss-120b.viol.pct}\% (gpt-oss-120b) to
\num{llama3.1-8b.viol.pct}\% (Llama 3.1 8B). Appeals to authority, deadlines and
claims that a rule had changed moved no model significantly in any of
\num{all.social.tests} tests. Counterfeit authorization did: a claimed compliance
sign-off with a reference number, or a forged system notice. The models did not all fall to the same counterfeit, so a single-attack
evaluation would have ranked them differently, and under the forged notice all eight, including the most compliant,
skipped a required approval. The prediction held on pressure wherever the new models could test it, and
failed on legitimate requests for one model, whose tool calls were written
into prose and never executed. Most violations were the agent's first action, and
in \num{all.misreport.k} episodes an agent denied an order the ledger shows it
placed. Since every violation was decidable before execution, these rules belong
in deterministic pre-trade controls rather than in the agent's judgment. Code,
scenarios and every transcript are released.
